% TOP_PROBLEM: {"id": 4600030, "title": "Decidability problems in symbolic dynamics", "category_id": 11, "category": {"id": 11, "name": "dynamical_systems", "display_name": "Dynamical Systems", "description": "Problems about long-term behavior of deterministic systems, Hamiltonian dynamics, and periodic orbits.", "slug": "dynamical-systems", "order_index": 11, "created_at": "2026-07-31T15:26:25.671Z"}, "status": "partially_solved", "problem_number": "AMR-045-0030", "set_id": 13, "statusQualification": "The right-inverse subquestion is decidable; only the remaining conjugacy and expansive-component questions are posed as research targets."}
% FIRST_POSTED: 2026-10-02
% ENTRY_KIND: statement_only
% UnsolvedMath ID 4600030; source catalogue code AMR-045-0030
% !TeX program = lualatex
% Definition and sources only; this file is not an LLM solution attempt.
\documentclass[11pt,a4paper]{article}
\usepackage[margin=25mm]{geometry}
\usepackage{fontspec}
\setmainfont{lmroman10-regular.otf}[BoldFont=lmroman10-bold.otf,
 ItalicFont=lmroman10-italic.otf,BoldItalicFont=lmroman10-bolditalic.otf]
\usepackage{amsmath,amssymb,mathtools}
\usepackage{unicode-math}
\setmathfont{latinmodern-math.otf}
\AtBeginDocument{\let\setminus\smallsetminus}
\usepackage{xurl,hyperref}
\hypersetup{colorlinks=true,urlcolor=blue}
\setcounter{secnumdepth}{0}
\setlength{\parindent}{0pt}
\setlength{\parskip}{5pt}
\setlength{\emergencystretch}{3em}
\begin{document}
\section*{Decidability problems in symbolic dynamics}\label{4600030}
{\small UnsolvedMath ID 4600030 · AMR-045-0030 · Dynamical Systems · AMR Open Problem Lists}
\subsection{Definitions and mathematical statement}
A shift of finite type (SFT) is a space of symbol sequences specified by
finitely many forbidden words. A sofic shift is the image of an SFT under
a continuous shift-commuting map, equivalently the label sequences along
paths in a finite labeled directed graph. Use either two-sided sequences
indexed by $\mathbb Z$ or one-sided sequences indexed by $\mathbb N$, as
specified. A conjugacy is a homeomorphism commuting with the shifts.

The principal decision questions ask for algorithms which halt on every
input and decide conjugacy of two given SFTs, two given two-sided sofic
shifts, or two given one-sided sofic shifts. Inputs are finite graph
presentations; no bound on a possible conjugating block rule is given.

A further question starts with an SFT shift $S$ and an expansive
automorphism $U$, defining $\alpha^{(a,b)}=S^aU^b$. A line $L\subset
\mathbb R^2$ is expansive if some $R,\varepsilon>0$ make observation of
$\alpha^v$ at every integer vector within distance $R$ of $L$ distinguish
any two distinct points by at least $\varepsilon$. Determine whether the
open component of expansive directions containing the $S$-direction can
be effectively approximated from both sides.

The catalogue also mentions deciding whether a surjective full-shift
block map has a continuous shift-commuting right inverse. That subquestion
is decidable by the split-epicness theorem cited below. It is retained
as resolved context; the conjugacy and expansive-direction questions
are the targets here.
\subsection{Short English statement}
Resolve the remaining algorithmic classification and expansive-direction questions for finite symbolic presentations.
\subsection{Sources}
\begin{itemize}
\item[S1] ULAM.AI and the UnsolvedMath Contributors, supplied problems.json, record 4600030 (AMR-045-0030), AMR Open Problem Lists. Catalogue identity and source transcription; CC BY 4.0.
\url{https://huggingface.co/datasets/ulamai/UnsolvedMath}.
\item[S2] Boyle - Open Problems in Symbolic Dynamics (2008), Problem/Question/Conjecture 19.1. Original source indexed by AMR-045-0030.
\url{https://www.math.umd.edu/~mboyle/papers/openfinalsub3nov2007.pdf}.
\item[S3] Salo, On von Neumann regularity of cellular automata, Theorem7; decidability of split epicness, following Salo–Törmä.
\url{https://arxiv.org/abs/2209.13373}.
\end{itemize}
\end{document}
